% !TeX root = ../../Book2.tex
% 纲要标题：### A.3 精确代数结论与近似输入的区别
% 纲要位置：计划纲要.md，第 1764 行。

% 纲要要点（供目录核对）：
%
% * 精确结构分类与近似输入所能支持的结论；区别输入区间不确定性与固定输入的舍入误差，有限二进制浮点数本身与其所代表的目标实数分别说明
% <!-- 短节：不设小节 -->
%

\section{精确代数结论与近似输入的区别}
\label{b2:sec:A03}

精确代数计算处理指定环或域中的恒等式、整除与相等。与近似计算比较时，须区分输入本身的不确定性和算法的舍入误差：只知道某个参数所在的区间，与已经固定一个输入、但运算过程中发生舍入，是两个不同问题。本节主要讨论前者。例如对实参数矩阵
\[
A_\varepsilon=\begin{pmatrix}1&1\\1&1+\varepsilon\end{pmatrix},
\]
精确行列式为 \(\varepsilon\)。因此 \(\varepsilon=0\) 时秩为一，而任意非零 \(\varepsilon\) 都给出秩二。如果输入只知道 \(|\varepsilon|\le10^{-12}\)，其中既包含零，也包含非零数；增加后续计算的精度不能消除输入本身留下的这一不确定性。

\ProseParagraphBreak
同样，矩阵
\[
N=\begin{pmatrix}0&1\\0&0\end{pmatrix},\qquad
N_\varepsilon=\begin{pmatrix}0&1\\\varepsilon&0\end{pmatrix}
\]
满足 \(N_\varepsilon\to N\)（\(\varepsilon\to0^+\)）。\(N\) 的最小多项式是 \(t^2\)，有一个二阶 Jordan 块；当实数 \(\varepsilon>0\) 时，\(N_\varepsilon\) 有两个不同的实特征值 \(\sqrt\varepsilon,-\sqrt\varepsilon\)，因此在实数域上可对角化，其 Jordan 形由这两个一阶块组成。Jordan 型中的块大小与数量是精确相似分类的离散数据；这个例子表明，任意小的条目扰动都可能改变它们。因而不能凭近似特征值是否相等来判断精确块结构，而须使用第四章的不变因子或核维数论证。

\ProseParagraphBreak
这也不表示精确算法能够从测量值恢复唯一的真实模型。一个有限的二进制浮点数本身是确定的数，可以精确解释为某个整数乘以二的整数次幂；近似性来自用这个数代表目标实数，浮点运算还可能引入新的舍入。把十进制字符串 \(0.1\) 指定为有理数 \(1/10\)，是在选择一个精确输入；通常的有限二进制浮点存储不能恰好表示 \(1/10\)，保存的是另一个二进有理数。若把该字符串解释为区间 \([0.0999,0.1001]\) 中的未知实数，输入又不再是单个确定值。显示相同的小数不保证这几种解释指定了同一个数学对象，题目应先明确采用哪一种。

\ProseParagraphBreak
在本书的代数算例中，整数以整数运算处理，有理数以分子分母处理，有限域元素按明确给出的模关系处理。输出应能通过精确乘法、余式为零、换基行列式在工作环中为单位等条件复核。数值计算则可以用于发现候选结构、估计解或提示某个矩阵接近降秩；若要把这些观察提升为精确命题，还须另给足够的代数证据。本节不讨论误差传播或数值稳定性的算法分析，这些问题属于数值线性代数的范围。
